\begin{python0}
from solutions import *; clear()
\end{python0}
\sectionthree{No constructor found}

Now we'll look at a very common case that you may run
into: the case when your compiler cannot find a constructor to call. For
example,

\begin{console}
class C
{
public:
        C(int a)
              : x_(a)
        {
              std::cout << "C()\ n";
        }
        
private:
        int x_;
};

class D
{
private:
        C c_;
};
\end{console}

\textbf{Error:} no appropriate constructor. Why? Examine the above code
carefully. Notice that the class \texttt{C} contains a constructor that
takes an \texttt{int} parameter.

So what's the problem? Well, if an object of type
\texttt{D} were created, then it would have to create its member,
\texttt{c}. However, an \texttt{int} parameter for the constructor is
necessary to create an object of type \texttt{C} (since \texttt{C} does not
have a default constructor). The question is, assuming this code
actually works, where would the \texttt{int} parameter come from? This
produces an error because there is \EMPHASIZE{no} \texttt{int} parameter
passed into the constructor.

No surprise there. Recall that you must pass in all parameters for any
function that takes them.

